\documentclass{article}
\usepackage{spconf,amsmath,amssymb,graphicx}
\usepackage[flushleft]{threeparttable}
\usepackage{cite}
\usepackage{booktabs,multirow}
\usepackage[table]{xcolor}
\usepackage[hidelinks]{hyperref}

\title{CASM: A Context-Aware Semi-Markov Alternative to DBN Post-Processing for Beat Tracking}

\name{Zhanhong He$^{1}$, Hanyu Meng$^{2}$, Yaolong Ju$^{3}$}
\address{
    $^{1}$The University of Western Australia, Perth, Australia\\
    $^{2}$The University of New South Wales, Sydney, Australia\\
    $^{3}$Great Bay University, Guangdong, China
}

\begin{document}
\ninept
\maketitle

\begin{abstract}
The last stage of beat tracking must turn dense neural activations into a
discrete and musically coherent event sequence. Direct peak picking preserves
local evidence but can fragment the pulse, whereas dynamic Bayesian networks
(DBNs) restore continuity through globally configured tempo and meter
assumptions that need not transfer to difficult music. We introduce CASM, a
plug-in, context-aware semi-Markov decoder over a sparse graph of
activation-supported beat candidates. CASM estimates local periodicity and its
ambiguity from the current activation sequence; this ambiguity jointly controls
the strength and width of each event-to-event duration potential. It therefore
commits strongly when the pulse is clear and relaxes toward the neural evidence
when competing metrical levels are plausible. A beat-count safeguard and a
beat-synchronous meter decoder make this behaviour deterministic and
auditable. One globally calibrated CASM configuration is applied without
backbone retraining or track-specific BPM and meter settings to Beat This,
MSCNN-lite, and TCN activations on GTZAN and SMC. The results position CASM as
a conservative F1--continuity operating point rather than a universal
replacement for DBN. An exploratory calibration-scale study further shows
metric-dependent saturation and decreasing sensitivity to fold composition.
\end{abstract}

\begin{keywords}
beat tracking, downbeat tracking, semi-Markov decoding, post-processing,
local periodicity
\end{keywords}

\section{Introduction}
\label{sec:intro}

The final decoder in a beat tracker is deceptively consequential. A neural
network can assign high framewise probability to genuine beats while still
producing duplicates, gaps, metrical-level switches, or downbeats that are
inconsistent with its beat grid. Event-level F-measure rewards accurate local
placement, whereas continuity measures additionally ask whether one metrical
interpretation persists through time. This creates a practical tension:
weak decoding respects a capable network but may yield an unstable clock;
strong decoding restores a clock but can overwrite correct, expressive, or
unusual evidence.

One prevalent structured decoder is the dynamic Bayesian network (DBN)
\cite{krebs2015dbn}. Its latent tempo, position, and meter states are valuable
precisely because they can repair noisy activations. The permitted tempo
support, meter set, and transition stiffness, however, are configured globally.
The posterior adapts its state path to the input, but the form and strength of
the temporal assumptions do not become more cautious when the input itself
reveals half/double-tempo ambiguity. Re-optimizing those settings on a target
corpus is legitimate development-set model selection, yet a
test- or target-informed recipe is no longer evidence of out-of-the-box
transfer. Recent post-processing-free trackers show that strong networks can
avoid some DBN errors \cite{chen2022postprocessingfree,foscarin2024beat},
but minimal peak picking can sacrifice sequence continuity. Recent failure
analysis on the deliberately difficult SMC corpus makes this trade-off
especially visible \cite{ahn2026smcblindspot}.

We ask whether the temporal prior itself can listen to the current activation.
Our answer is \emph{CASM}, a Context-Aware Semi-Markov post-processor. CASM
constructs a sparse graph whose nodes are activation maxima and whose edges are
complete inter-beat segments. At both endpoints of an edge, it estimates a
local period and a best-versus-competing-period margin. Clear periodic evidence
creates a narrow, strong duration potential; ambiguous evidence simultaneously
weakens and broadens that potential. CASM is consequently an automatic
correction layer in a precise, limited sense: after one global calibration, it
requires neither track metadata nor per-track decoder adjustment, while its
constraints remain input-dependent at inference.

Our contributions are threefold:
\begin{itemize}
  \setlength\itemsep{0pt}
  \setlength\parskip{0pt}
  \item We formulate maximum-score beat decoding on a sparse semi-Markov graph
  of activation-supported candidates. Each duration potential uses both
  endpoint contexts and is conditioned on their ambiguity.
  \item We connect beat decoding to a beat-synchronous variable-meter decoder
  and two explicit safeguards, yielding a deterministic plug-in that does not
  retrain the upstream tracker.
  \item We study one frozen operating point across three different activation
  backbones and analyse how the amount and composition of decoder-calibration
  data affect transfer, while separating exploratory selection from clean
  evaluation claims.
\end{itemize}

\section{Related Work}
\label{sec:related}

\subsection{From activations to event sequences}

Dynamic programming has long combined onset evidence with inter-beat
regularity \cite{ellis2007beattracking}. Korzeniowski \emph{et al.} proposed a
Bayesian beat-position extractor around a dominant beat interval
\cite{korzeniowski2014crf}; its implementation is exposed by \texttt{madmom}
under the historical name \texttt{CRFBeatDetectionProcessor}, although the
paper is not the learned linear-chain CRF of Fillon \emph{et al.}
\cite{fillon2015crf}. The learned CRF branch was extended to downbeats
\cite{durand2016crfdownbeat} and, in 2019, to skip-chain musical structure and
Afro-Latin microtiming \cite{fuentes2019skipchain,fuentes2019microtiming}.
Resonating comb filters form a different branch: they estimate a dominant
global tempo from recurrent activations \cite{bock2015combfilter}, after which
events are aligned to that periodicity. The 2019 multi-task tempo--beat model
\cite{bock2019multitask} extended neural training rather than this comb-filter
event decoder, and retained structured inference. DBNs instead enumerate
tempo--phase--meter states and decode their most likely trajectory
\cite{krebs2015dbn,bock2016joint}; they remain common in recent high-performing
systems.

\subsection{Local, segmental, and modern alternatives}

Reliability-informed tracking established the principle that temporal
commitment should depend on evidence quality \cite{degara2012reliability}.
More recently, a compact one-dimensional semi-Markov state space represented
time-varying rhythmic jumps for causal joint tracking
\cite{heydari2022semimarkov}. Particle filters have long maintained online
tempo hypotheses \cite{hainsworth2004particle}, and BeatNet/BeatNet+ extend
that branch to joint real-time rhythm analysis
\cite{heydari2021beatnet,heydari2024beatnetplus}.
PLPDP is CASM's closest conceptual predecessor: it derives predominant local
pulse information and uses dense-frame dynamic programming for expressive
tempo \cite{chiu2023localperiodicity}. Its real-time continuation improves
controllability \cite{meier2024realtime}, and dPLP makes local-pulse estimation
differentiable \cite{chiu2025dplp}.

CASM does not claim novelty for local periodicity, dynamic programming, or
semi-Markov inference separately. Unlike PLPDP's dense frame grid, it scores
only segments joining retained activation maxima and conditions each segment on
the local period \emph{and ambiguity at both endpoints}. Ambiguity modulates
both penalty strength and tolerance, rather than merely supplying a preferred
period. Unlike the one-dimensional causal semi-Markov model, CASM is an
offline sparse event graph with explicit fallback behaviour. Unlike a DBN, it
does not maintain a dense tempo--phase state lattice; unlike a semi-CRF
\cite{sarawagi2004semimarkov}, its potentials are not learned or normalized.

Another response to difficult rhythm is to adapt the tracker with user
annotations \cite{pinto2021userdriven,maia2022adapting,pinto2023bridging,
maia2024selective,pinto2026challenging}. That line is complementary: it deliberately imports
target-specific supervision. CASM instead recomputes its constraint from
unlabelled test activations and leaves upstream weights fixed. At the other
extreme, recent work reformulates the prediction task itself
\cite{bolt2026reformulated}; masked diffusion learns coherent sequences through
iterative refinement \cite{foscarin2026masked}, while BeatFCOS predicts
intervals as objects with Soft-NMS \cite{ahn2025beatfcos}. Knowledge-driven
metrical hypotheses have also been introduced during self-supervised
pre-training \cite{gagnere2025multiple}. These are modern alternatives to
traditional post-processing, but require a changed or retrained predictor;
CASM targets the still-useful deployment setting in which only a frozen
activation interface is available.

\section{CASM: Context-Aware Semi-Markov Decoding}
\label{sec:method}

\subsection{Candidate graph and segmental objective}

Let a frozen tracker output beat and downbeat logits
$z_t^{\mathrm b}$ and $z_t^{\mathrm d}$ at $F$ frames per second, and let
$a_t=\operatorname{sigmoid}(z_t^{\mathrm b})$. We retain one-dimensional
local maxima, including valid boundary maxima, whose probability exceeds
$\theta_{\mathrm c}$. The ordered candidate set is
$\mathcal C=\{t_1,\ldots,t_N\}$. A directed edge $(i,j)$ spans the complete
candidate-to-candidate interval
\[
  \Delta_{ij}=(t_j-t_i)/F
\]
and is admissible when
$\Delta_{ij}\in[60/B_{\max},60/B_{\min}]$. CASM never inserts an event away
from a retained neural maximum; it selects a sparse path
$\pi=(i_1,\ldots,i_K)$ by maximizing
\begin{equation}
  \mathcal S(\pi)=\sum_{k=1}^{K}e_{i_k}
  -\sum_{k=2}^{K}D_{i_{k-1},i_k}.
  \label{eq:casm-objective}
\end{equation}
Here $e_i$ is clipped logit evidence and $D_{ij}$ scores the duration of a
whole segment. This variable-duration edge potential is the precise
semi-Markov component; Eq.~\eqref{eq:casm-objective} is a maximum-score
decoder, not a generative hidden semi-Markov model.

\subsection{Local period and ambiguity}

For each candidate $i$, we evaluate integer lags
$\ell\in\mathcal L=\{\ell_{\min},\ldots,\ell_{\max}\}$ in a centred local
window $\mathcal W_i$. Zero-padded normalized autocorrelation gives
\begin{equation}
q_i(\ell)=
\frac{\langle a_u a_{u-\ell}\rangle_{u\in\mathcal W_i}}
{\sqrt{\langle a_u^2\rangle_{u\in\mathcal W_i}
       \langle a_{u-\ell}^2\rangle_{u\in\mathcal W_i}+\epsilon_q}}
\left(\frac{\ell_{\min}}{\ell}\right)^\beta ,
\label{eq:casm-period}
\end{equation}
where the last factor is a mild fast-tempo bias. We take
$\ell_i^\star=\arg\max_\ell q_i(\ell)$ and $\tau_i=\ell_i^\star/F$.
The runner-up excludes the best lag and its two neighbours, exposing
competing nearby and metrical-level periodicities. Their separation is
\begin{equation}
c_i=\operatorname{clip}_{[0,1]}
\left[
\frac{q_i(\ell_i^\star)-
\max_{\ell:\,|\ell-\ell_i^\star|>2}q_i(\ell)}
{|q_i(\ell_i^\star)|+\epsilon_c}
\right].
\label{eq:casm-ambiguity}
\end{equation}
We call $c_i$ a \emph{period margin}, not calibrated correctness. If at least
five candidates exist, both $\tau_i$ and $c_i$ are median-filtered over five
successive candidates.

\subsection{Ambiguity-conditioned duration inference}

For edge $(i,j)$, CASM combines the endpoint contexts geometrically:
$\bar\tau_{ij}=\sqrt{\tau_i\tau_j}$ and
$c_{ij}=\sqrt{c_i c_j}$. Its duration cost is
\begin{equation}
D_{ij}=\frac{\lambda c_{ij}}{2}
\left[
\frac{\log(\Delta_{ij}/\bar\tau_{ij})}
{\sigma_0+(1-c_{ij})\sigma_{\mathrm u}}
\right]^2 .
\label{eq:casm-duration}
\end{equation}
The log ratio treats proportional timing errors symmetrically. More
importantly, $c_{ij}$ plays two coupled roles. A separated local-period peak
increases the multiplier and narrows the scale, enforcing a clear pulse.
Competing periods reduce the multiplier and widen the scale, allowing the path
to follow activation evidence instead of confidently imposing the wrong
metrical level.

With
$e_j=\operatorname{clip}(z^{\mathrm b}_{t_j}/T,-L,L)$, the Viterbi-style
recurrence is
\begin{equation}
V_j=e_j+\max\!\left\{0,\!
\max_{\substack{i<j\\(i,j)\ {\rm admissible}}}
\left(V_i-D_{ij}\right)\right\}.
\label{eq:casm-dp}
\end{equation}
The zero branch permits a new path after an unfavourable prefix, and
backtracking starts from $\arg\max_j V_j$. In parallel, a Direct path is
formed by seven-frame max pooling, a strict zero-logit threshold, and
one-frame peak deduplication, matching Beat This's minimal decoder
\cite{foscarin2024beat}. If the number of structured events divided by the
number of Direct events lies outside
$[r_{\min},r_{\max}]$, the Direct path is returned. Empty candidate input
returns an empty sequence. The count test is a deterministic safety gate, not
a learned router.

\subsection{Beat-synchronous downbeats}

Let the decoded beat grid be
$\widehat{\mathcal B}=\{\widehat b_n\}$. CASM scores possible downbeats only
at these beat locations. For a bar length $m$ in the allowed meter set
$\mathcal M$, define
$g_n=\operatorname{clip}(z^{\mathrm d}_{\widehat b_n}/T_{\mathrm d},-L,L)
+\rho$. A variable-meter recurrence connects successive downbeats:
\begin{equation}
H_{n,m}=g_n+\max_{m'\in\mathcal M}
\left[H_{n-m,m'}-\eta\,\mathbf 1(m'\ne m)\right].
\label{eq:casm-meter}
\end{equation}
Initial states are permitted near the beginning and termination near the end,
so a decoded path can change bar length while paying switch cost $\eta$.
Separately, Direct downbeat maxima are snapped to
$\widehat{\mathcal B}$. The meter path is retained only if its one-to-one
event F-measure against this snapped path is at least $\gamma$ under a 70-ms
tolerance; otherwise CASM returns the snapped path. Consequently every
reported downbeat remains on the selected beat grid.

\section{Experiments}
\label{sec:exp}

\subsection{Data, backbones, and metrics}

We evaluate on GTZAN ($N=993$) using the annotations and exclusions of Beat
This \cite{foscarin2024beat}, and on the deliberately difficult SMC corpus
($N=217$) \cite{smc}; SMC provides beat but not downbeat annotations. We use
three 50-Hz activation backbones: the Beat This spectrogram Transformer
\cite{foscarin2024beat}, the lightweight multi-scale convolutional backbone
MSCNN-lite \cite{he2026dynest}, and the TCN of
\cite{bock2020deconstruct}. The decoder receives only their framewise beat
and downbeat logits.

For GTZAN, \texttt{final0}, \texttt{final1}, and \texttt{final2} are three
independently trained full-data backbone seeds. Their training excludes
GTZAN but includes SMC. We first macro-average over the same 993 tracks and
then average the three seed-level scores. For SMC, using those full-data
checkpoints would leak training examples; we instead use the corresponding
held-out checkpoint for each of eight folds. Seven folds contain 27 tracks and
the eighth contains 28; we average their fold-level macro scores. CASM itself
is deterministic and has neither a neural checkpoint nor a random seed.

Following the established protocol, the first 5\,s are trimmed and
\texttt{mir\_eval} computes F1 with $\pm70$\,ms tolerance, CMLt, and AMLt
\cite{raffel2014mireval}. CMLt requires a continuously correct metrical
level; AMLt additionally accepts the standard off-beat and half/double-level
interpretations. Seed variation on GTZAN and fold variation on SMC have
different replication units and are descriptive, not confidence intervals.

\subsection{Decoder baselines}

\emph{Direct} is the minimal peak picker defined in
Sec.~\ref{sec:method}. \emph{DBN} is the conventional joint
\texttt{madmom} configuration: meters $\{3,4\}$, 55--215 BPM,
transition weight 100, observation weight 16, and threshold 0.05
\cite{krebs2015dbn}. \emph{DBN-SMC-opt} is deliberately diagnostic. Its
sequential beat and bar stages were selected from 753 configurations on the
TCN SMC 8-fold OOF activations, maximizing beat F1 with CMLt/AMLt tie-breaks,
and only then frozen for the three backbones and GTZAN. It uses 50--215 BPM,
30 tempo states, transition/observation weights 7/20, threshold 0.3,
meters $\{3,4\}$, bar-observation weight 100, and meter-switch probability
$10^{-7}$. It is neither a clean SMC test nor a published
SMC-specific recipe from \cite{bock2020deconstruct}.

The \emph{CRF} row uses \texttt{madmom}'s historically named
\texttt{CRFBeatDetectionProcessor} \cite{korzeniowski2014crf} at 30--300 BPM
with interval width 0.18; it should not be confused with the learned CRF in
\cite{fillon2015crf}. \emph{CombFilter} uses \texttt{madmom}'s global
tempo/iterative-alignment processor \cite{bock2015combfilter}. CRF and
CombFilter decode downbeat activations at 8--100 BPM and snap their outputs to
their respective beat grids.

\subsection{Frozen operating point and scale analysis}

Every CASM row in Sec.~\ref{sec:results} uses exactly one explicitly loaded
Frozen-4F configuration:
\[
\begin{split}
&F=50,\quad \theta_{\mathrm c}=0.03,\quad
(B_{\min},B_{\max})=(30,300),\\
&|\mathcal W|=8{\rm\,s},\quad \beta=0.15,\quad
(T,L)=(2,6),\\
&(\lambda,\sigma_0,\sigma_{\mathrm u})=(4,0.12,0.40),\quad
[r_{\min},r_{\max}]=[0.85,1.8],\\
&\mathcal M=\{2,\ldots,7\},\quad
(T_{\mathrm d},\eta,\rho,\gamma)=(4,3,0,0.6).
\end{split}
\]
Its scalar settings were selected on the union of SMC folds
$\{2,7,1,3\}$ and then reused for all tracks and all backbones. Thus
context-aware means input-conditioned inference, not test-time parameter
updates: CASM has validation-selected parameters, but no annotated beat,
inter-beat-interval, BPM, or meter input at inference.

The choice to foreground 4F was made \emph{after} inspecting the calibration
sensitivity analysis and is therefore exploratory. The seven-fold
configuration is the only comparator locked before GTZAN inspection; it
differs from 4F only in $\sigma_0=0.15$. For the scale study, SMC fold 0
($N=27$) is never used for selection. From folds 1--7 we enumerate every
1-, 2-, and 4-fold calibration subset ($7$, $21$, and $35$ combinations) plus
the seven-fold union, and evaluate all locked configurations on the same SMC
fold 0 and fixed GTZAN/\texttt{final1} panel. The latter panel is
test-conditioned and is used only to describe sensitivity, not to claim a
held-out optimum.

\section{Results and Discussion}
\label{sec:results}

\subsection{A single operating point across backbones}

% DATA AUDIT (2026-09-04): Numeric cells below are deliberately preserved at
% the owner's request. Reconcile them with the 2026-09-03 canonical absolute
% ledger before submission; do not infer new values from the displayed deltas.
\begin{table*}[t]
\centering
\caption{Post-processor comparison. Direct rows are absolute percentages;
all following rows are percentage-point changes from the Direct row of the
same backbone. CASM uses one Frozen-4F configuration. DBN-SMC-opt is a
target-selected diagnostic, not a clean SMC result.}
\label{tab:postprocessor_ablation}
\begingroup
\fontsize{7.5pt}{9pt}\selectfont
\setlength{\tabcolsep}{3.2pt}
\renewcommand{\arraystretch}{1.0}
\resizebox{\textwidth}{!}{%
\begin{tabular}{l|rrr|rrr|rrr}
\toprule
\multicolumn{1}{c|}{\multirow{2}{*}{\textbf{Method}}}
& \multicolumn{6}{c|}{\textbf{GTZAN (Beat / Downbeat)}}
& \multicolumn{3}{c}{\textbf{SMC (Beat)}} \\
\cmidrule(lr){2-7}
\cmidrule(l){8-10}
& F1 & CMLt & AMLt
& F1 & CMLt & AMLt
& F1 & CMLt & AMLt \\
\midrule

\textbf{BeatThis} (direct) \cite{foscarin2024beat}
& 89.1 & 79.8 & 89.8
& 78.3 & 67.2 & 79.1
& 62.7 & 51.4 & 61.0 \\
\rowcolor{gray!10}
\quad w. CASM
& $\mathbf{+0.2}$ & $+0.4$ & $+0.6$
& $\mathbf{+0.6}$ & $+3.2$ & $+6.0$
& $\mathbf{+0.3}$ & $\mathbf{+2.4}$ & $+3.3$ \\
\quad w. DBN
& $-1.0$ & $\mathbf{+0.7}$ & $+1.3$
& $-0.9$ & $\mathbf{+6.2}$ & $\mathbf{+8.7}$
& $-5.2$ & $-3.9$ & $+1.5$ \\
\rowcolor{gray!10}
\quad w. DBN, SMC-opt
& $-0.5$ & $-0.9$ & $+0.2$
& $-1.5$ & $-0.6$ & $+1.1$
& $-1.6$ & $+0.4$ & $\mathbf{+3.5}$ \\
\quad w. CombFilter \cite{bock2015combfilter}
& $-4.0$ & $-3.6$ & $+0.1$
& $-3.3$ & $+3.2$ & $+7.6$
& $-7.1$ & $+1.1$ & $+3.1$ \\
\rowcolor{gray!10}
\quad w. CRF \cite{korzeniowski2014crf}
& $-0.6$ & $-0.3$ & $\mathbf{+1.4}$
& $-1.4$ & $+1.9$ & $+5.5$
& $-1.2$ & $+1.4$ & $+1.2$ \\
\midrule
\textbf{MSCNN-lite} (direct) \cite{he2026dynest}
& 86.1 & 73.4 & 79.9
& 70.8 & 53.1 & 63.1
& 57.2 & 42.0 & 48.4 \\
\rowcolor{gray!10}
\quad w. CASM
& $\mathbf{+0.4}$ & $+1.5$ & $+2.4$
& $+1.7$ & $+6.8$ & $+11.7$
& $\mathbf{+0.5}$ & $+3.3$ & $+4.4$ \\
\quad w. DBN
& $-0.6$ & $\mathbf{+4.0}$ & $\mathbf{+7.2}$
& $\mathbf{+1.8}$ & $\mathbf{+15.8}$ & $\mathbf{+21.2}$
& $-4.9$ & $+0.3$ & $+3.9$ \\
\rowcolor{gray!10}
\quad w. DBN, SMC-opt
& $-1.1$ & $+0.2$ & $+5.9$
& $-0.9$ & $+9.5$ & $+14.6$
& $+0.3$ & $\mathbf{+3.5}$ & $\mathbf{+6.8}$ \\

\midrule
\textbf{TCN} (direct) \cite{bock2020deconstruct}
& 87.1 & 73.8 & 81.7
& 62.4 & 24.6 & 59.9
& 52.9 & 30.1 & 36.0 \\
\rowcolor{gray!10}
\quad w. CASM
& $\mathbf{+0.3}$ & $\mathbf{+1.3}$ & $+1.6$
& $+2.3$ & $+19.4$ & $+14.0$
& $+0.3$ & $+3.8$ & $+5.2$ \\
\quad w. DBN
& $-2.6$ & $-0.7$ & $\mathbf{+5.0}$
& $\mathbf{+4.9}$ & $\mathbf{+23.2}$ & $\mathbf{+22.9}$
& $-0.9$ & $-7.0$ & $+2.3$ \\
\rowcolor{gray!10}
\quad w. DBN, SMC-opt
& $-1.9$ & $-0.8$ & $+4.4$
& $-0.2$ & $+8.9$ & $+9.3$
& $\mathbf{+1.7}$ & $\mathbf{+4.5}$ & $\mathbf{+8.2}$ \\
% Completed TCN/MSCNN-lite CombFilter and CRF rows exist in the canonical
% ledger and should be integrated only when the owner unlocks table revision.
\bottomrule
\end{tabular}%
}
\endgroup
\end{table*}

Table~\ref{tab:postprocessor_ablation} is best read as a transfer and
trade-off table, not a winner count. At its displayed precision, the same
CASM setting produces positive changes from Direct across the three
backbones, two corpora, and all available metrics. Its characteristic
behaviour is conservative: it preserves local event placement while recovering
part of the continuity absent from peak picking. DBN produces larger
continuity changes in several GTZAN downbeat conditions, but those changes are
frequently accompanied by lower beat or downbeat F1 and are less stable on
SMC. The SMC-selected DBN also changes its trade-off when transferred back to
GTZAN. We therefore do not claim that CASM dominates DBN, CRF, or comb
filtering. The evidence supports a reusable operating point whose corrections
depend on the input rather than a new target-specific recipe.

\subsection{Context among recent trackers}

% DATA AUDIT (2026-09-04): Numeric cells below are deliberately preserved.
% Three cells conflict with the 2026-09-03 canonical ledger and must be
% adjudicated by the owner before submission.
\begin{table*}[t]
\centering
\caption{Contextual comparison with recent beat/downbeat trackers. Values are
percentages reported under each cited system's protocol; differing training
data, checkpoints, and inference procedures make this literature context, not
a controlled ranking. Our rows use GTZAN three-seed means and SMC eight-fold
OOF means.}
\label{tab:sota}
\begingroup
\fontsize{7.5pt}{9pt}\selectfont
\setlength{\tabcolsep}{2.6pt}
\renewcommand{\arraystretch}{1.0}
\resizebox{\textwidth}{!}{%
\begin{tabular}{l|ccc|ccc|ccc}
\toprule
\multicolumn{1}{c|}{\multirow{2}{*}{\textbf{Method}}}
& \multicolumn{6}{c|}{\textbf{GTZAN (Beat / Downbeat)}}
& \multicolumn{3}{c}{\textbf{SMC (Beat)}} \\
\cmidrule(lr){2-4}
\cmidrule(lr){5-7}
\cmidrule(l){8-10}
& F1 & CMLt & AMLt
& F1 & CMLt & AMLt
& F1 & CMLt & AMLt \\
\midrule
\rowcolor{gray!10}
\textbf{HN-MERT} w. DBN \cite{ru2025hingenet}
& \textbf{89.7} & 81.4 & \textbf{94.3}
& 77.4 & 71.6 & 87.2
& 61.7 & -- & -- \\
\textbf{HN-MusicFM} w. DBN
& 89.2 & 80.9 & 93.7
& \textbf{79.8} & 73.2 & \textbf{89.5}
& 60.6 & -- & -- \\
\rowcolor{gray!10}
\textbf{BeatFM-MERT} w. DBN \cite{ru2025beatfm}
& \underline{89.5} & \underline{82.7} & \underline{93.9}
& 76.7 & 70.2 & 85.4
& 61.3 & -- & -- \\
\textbf{BeatFM-MusicFM} w. DBN
& 89.1 & 80.6 & 93.5
& \underline{79.6} & \underline{74.4} & \underline{88.7}
& 60.5 & -- & -- \\
\rowcolor{gray!10}
\textbf{BeatMamba} w. DBN \cite{ru2026beatmamba}
& 88.9 & 82.3 & \textbf{94.3}
& -- & -- & --
& 58.7 & 46.8 & 62.4 \\
\textbf{BeatKAN} w. DBN \cite{zhang2025beatkan}
& 88.2 & 78.1 & 92.3
& -- & -- & --
& 59.8 & 48.4 & \underline{64.0} \\
\rowcolor{gray!10}
\textbf{MDM-BeatThis} (direct) \cite{foscarin2026masked}
& \textbf{89.7} & \textbf{82.9} & 92.5
& 79.5 & \textbf{76.4} & 88.5
& 61.0 & 48.9 & \textbf{64.3} \\
\textbf{BeatThis} (direct) \cite{foscarin2024beat}
& 89.1 & 79.8 & 89.8
& 78.3 & 67.2 & 79.1
& \underline{62.7} & \underline{51.4} & 61.0 \\
\rowcolor{gray!10}
\textbf{BeatThis} w. DBN
& 88.1 & 80.5 & 91.1
& 77.4 & 73.4 & 87.8
& 57.5 & 47.5 & 62.5 \\
\textbf{BeatThis} w. CASM (proposed)
& 89.3 & 80.2 & 90.4
& 78.9 & 70.4 & 85.1
& \textbf{63.0} & \textbf{53.8} & \textbf{64.3} \\
\bottomrule
\end{tabular}%
}
\vspace{2pt}
\begin{minipage}{0.99\textwidth}
\footnotesize
GTZAN rows from this work average three independently trained backbone seeds;
SMC rows from this work average eight fold-level macro scores from 217
track-disjoint OOF predictions. These are different descriptive replication
axes. CASM uses the same Frozen-4F decoder in both cases; its SMC row is
backbone-level OOF but not strict nested OOF for decoder selection.
\end{minipage}
\endgroup
\end{table*}

Table~\ref{tab:sota} answers a narrower question than
Table~\ref{tab:postprocessor_ablation}: whether an accessible modern backbone
remains competitive when its minimal decoder is replaced. We could reproduce
a compatible activation interface only for Beat This, so the table does not
imply that CASM improves HingeNet, BeatFM, BeatMamba, BeatKAN, or MDM. Their
rows retain their authors' training and decoding pipelines. This distinction
also clarifies CASM's practical niche. MDM and BeatFCOS alter the learned
predictor to produce coherent sequences or intervals
\cite{foscarin2026masked,ahn2025beatfcos}; CASM can instead be attached to
already trained logit-producing models. Testing the same decoder on future
systems requires authors to expose framewise activations, not only final event
times.

\subsection{What calibration-scale sensitivity shows}

% The full nine-panel plot remains in figures/p2.jpg for a supplement.  In the
% main paper, its size displaces the references and obscures the central result.
The fixed-panel sensitivity analysis does not show that 4F is optimal, nor
that more data hurts machine learning. CASM learns no network weights here:
the independent variable is the amount of data used to select global decoder
scalars. On the held-out SMC panel, the family centres of F1, CMLt, and AMLt
rise from 1F to 7F. On GTZAN, beat F1 and downbeat F1/CMLt peak at 2F, while
beat CMLt/AMLt and downbeat AMLt continue to 7F. Hence the defensible
conclusion is metric-dependent diminishing returns rather than a universal
best scale.

There is nevertheless a useful stability result. From 1F to 2F to 4F, the
across-combination population spread decreases in all nine panels: more
calibration data makes the selected decoder less sensitive to which SMC folds
are present, even when the external point estimate does not improve
monotonically. We retain 4F as a balanced, post-hoc operating point for the
present analysis. Its near-identical seven-fold comparator is the appropriate
pre-GTZAN locked reference; a new external test set or nested selection is
needed to establish superiority.

\subsection{Limitations and decisive next tests}

Three limitations bound the current claim. First, SMC activations are
backbone-level OOF, but the single 4F decoder is not selected in a nested outer
loop. Second, CASM searches 30--300 BPM while the default DBN searches
55--215 BPM. Recent SMC analysis reports many references below that lower
limit \cite{ahn2026smcblindspot}; a matched-support DBN is therefore required
before attributing an SMC difference solely to CASM's ambiguity mechanism.
Third, the present aggregate tables do not isolate the sparse graph,
strength-versus-width ambiguity gate, or the two safeguards. A direct PLPDP
baseline, fixed-period DP, component ablations, and per-corpus fallback and
meter-path acceptance rates are the decisive experiments. Track-paired
bootstrap or randomization tests should accompany them; three training seeds
or eight folds alone do not justify significance claims.

\section{Conclusion and Future Work}
\label{sec:concl}

CASM revisits structured beat decoding as an inductive bias conditioned on each
input, not as a return to a fixed global grid. It joins activation-supported events
with locally centred segment potentials and deliberately relaxes those
potentials when periodic evidence is ambiguous. Once globally calibrated,
the same deterministic module runs without upstream retraining, user
annotation, or per-track tempo and meter settings. Current evidence supports
a conservative F1--continuity alternative to Direct and DBN decoding, not
universal dominance. The next version should use strict nested or
leave-one-dataset-out calibration, matched tempo-support baselines, PLPDP and
mechanism ablations, and explicit safeguard statistics. Multi-hypothesis
period paths, causal decoding, and a differentiable duration layer are natural
extensions, as is evaluation on every recent tracker that releases compatible
activations.

\bibliographystyle{IEEEbib}
\bibliography{casm}

\end{document}
